\par\medskip\noindent\textbf{(i) Vollständige Prüfung der assoziativen Tafeln.}\par
Jede zweielementige Halbgruppe wird durch eine Nummerierung
\(e\colon S_2\bij A\) auf eine assoziative Operation auf
\(S_2=\{0,1\}\) transportiert.  Nach
\FormulaRefAuto{FiniteSemigroupRelabelOrbitsAreIsomorphismClasses}[thm]
entsprechen deren Umbenennungsorbits genau den Isomorphieklassen.
Wir schreiben eine Tafel auf \(S_2\) zeilenweise als
\[
  t=(t_{00},t_{01},t_{10},t_{11})\in\{0,1\}^{4}.
\]
Genau die folgenden acht Tupel sind assoziativ:
\[
\begin{array}{cccc}
0000&0001&0011&0101\\
0110&0111&1001&1111.
\end{array}
\]
Für die konstanten Operationen sind beide Klammerungen gleich dem konstanten
Wert; für die linke beziehungsweise rechte Projektion sind sie beide gleich
\(x\) beziehungsweise \(z\).  Ferner gilt
\[
  \min(\min(x,y),z)=\min(x,y,z)=\min(x,\min(y,z)),
\]
und entsprechend für das Maximum.  Das Tupel \(0110\) beschreibt
\(x\star y\equiv x+y\pmod2\), das Tupel \(1001\) beschreibt
\(x\star y\equiv x+y+1\pmod2\).  In beiden Fällen ergeben die zwei
Klammerungen denselben Rest modulo \(2\).
Für jedes der übrigen acht Tupel liefert das angegebene Tripel verschiedene
Werte der beiden Klammerungen:
\[
\begin{array}{c|c}
\text{Tafeltupel}&\text{Tripel }(x,y,z)\text{ mit }(xy)z\neq x(yz)\\
\hline
0010,\ 0100&(1,0,1)\\
1000,\ 1110&(0,0,1)\\
1010,\ 1011,\ 1100,\ 1101&(0,0,0).
\end{array}
\]
Damit sind alle \(2^4=16\) Tafeln erfasst.

\par\medskip\noindent\textbf{(ii) Bestimmung der fünf Isomorphieklassen.}\par
Die nichttriviale Umbenennung \(0\leftrightarrow1\) zerlegt die acht
assoziativen Tafeln in genau die fünf Orbits
\[
\{0000,1111\},\quad
\{0001,0111\},\quad
\{0011\},\quad
\{0101\},\quad
\{0110,1001\}.
\]
Sie entsprechen der konstanten Halbgruppe, dem Zweierhalbverband, der
Linksnullhalbgruppe, der Rechtsnullhalbgruppe und der Zweiergruppe.

\par\medskip\noindent\textbf{(iii) Unterscheidung durch die Potenzhalbgruppen.}\par
Alle zugehörigen Potenzhalbgruppen besitzen drei Elemente.  Die folgenden
Eigenschaften hängen nur von ihrer abstrakten Multiplikation ab:
\begin{center}
\small
\begin{tabularx}{\linewidth}{@{}>{\raggedright\arraybackslash}Xc
  >{\raggedright\arraybackslash}X@{}}
\toprule
\textbf{Grundtyp und Regel}&
\textbf{Idempotente in \(\PowNE{S}\)}&
\textbf{Kennzeichen der Potenzhalbgruppe}\\
\midrule
Nullhalbgruppe (\(xy=0\))&1&konstant und kommutativ\\
Linksnullhalbgruppe (\(xy=x\))&3&\(XY=X\), nicht kommutativ\\
Rechtsnullhalbgruppe (\(xy=y\))&3&\(XY=Y\), nicht kommutativ\\
Zweierhalbverband (\(xy=\max\{x,y\}\))&3&kommutativ\\
Zweiergruppe (\(C_2\))&2&kommutativ\\
\bottomrule
\end{tabularx}
\end{center}
Die Anzahl der Idempotenten trennt zunächst die konstante Halbgruppe und die
Zweiergruppe ab.  Von den drei übrigen Fällen ist nur der Halbverband
kommutativ.  Links- und Rechtsnullfall sind ebenfalls getrennt: Wäre
\(F\) ein Isomorphismus vom Linksnull- zum Rechtsnullfall, so gälte für
\(X\neq Y\)
\[
  F(X)=F(XY)=F(X)F(Y)=F(Y),
\]
im Widerspruch zur Injektivität.  Isomorphe Potenzhalbgruppen gehören daher
zum selben der fünf Grundtypen, und die Ausgangshalbgruppen sind isomorph.

Die folgenden Schlusszeilen verwenden die vollständige Fallliste (i), ihre
Umbenennungsorbits (ii) und die multiplikativ definierten Kennzeichen (iii).
Für (iii) genügt Produkterhaltung in beiden Richtungen: Ein Element ist
idempotent genau dann, wenn sein Bild idempotent ist; Kommutativität wird
wegen der Surjektivität übertragen. Die Links- beziehungsweise Rechtsnullregel
wird ebenfalls für jedes Paar übertragen. Damit sind sämtliche verwendeten
Kennzeichen Isomorphieinvarianten.
\begin{tabproofwide}
  \proofstepwidestar[1]{\FiniteSemigroupAlg{A}{\star}}{\rA}
  \proofstepwidestar[2]{\FiniteSemigroupAlg{B}{\diamond}}{\rA}
  \proofstepwidestar[3]{\FiniteCard(A)=2}{\rA}
  \proofstepwidestar[4]{\FiniteCard(B)=2}{\rA}
  \proofstepwidestarbreak[5]{\SemigroupIso{F}
    {\PowNE{A},\star_{\mathcal P}}
    {\PowNE{B},\diamond_{\mathcal P}}}{\rA}
  \proofstepwidestar[1]{\SemigroupAlg{A}{\star}}{%
    \FormulaRefAuto{FiniteSemigroupBaseAxiom}[ax]{1}}
  \proofstepwidestar[2]{\SemigroupAlg{B}{\diamond}}{%
    \FormulaRefAuto{FiniteSemigroupBaseAxiom}[ax]{2}}
  \proofstepwidestarbreak[1,2,3,4,5]{\operatorname{can}(A,\star)
    =\operatorname{can}(B,\diamond)}{%
    \textnormal{\parbox[t]{\linewidth}{\raggedright
      Fallliste (i), Orbits (ii), Invarianten (iii)}}}
  \proofstepwidestarbreak[1,2]{\begin{aligned}[t]
    &\bigl(\exists f\,\FiniteSemigroupIso{f}{A,\star}{B,\diamond}\bigr)\\[-2pt]
    &\quad\leftrightarrow\operatorname{can}(A,\star)=\operatorname{can}(B,\diamond)
    \end{aligned}}{%
    \FormulaRefAuto{FiniteSemigroupCanonicalTableCodeComplete}[thm]{1,2}}
  \proofstepwidestarforcebreak[1,2,3,4,5]{\exists f\,
    \FiniteSemigroupIso{f}{A,\star}{B,\diamond}}{\rRE{\rLREb{9},8}}
\end{tabproofwide}
